% authority: science-draft
% status: draft/control
% route: ontology-law-research-packet
% trigger_classification: derivation_critical_missing_source_law
% adoption_status: blocked_adoption_open_continuation
\documentclass[11pt]{article}
\usepackage[margin=1in]{geometry}
\usepackage{amsmath,amssymb,amsthm}
\usepackage{booktabs,longtable,array}
\usepackage[T1]{fontenc}
\usepackage{enumitem}

\newtheorem{definition}{Definition}
\newtheorem{theorem}{Theorem}
\newtheorem{corollary}{Corollary}
\newtheorem{proposition}{Proposition}
\newtheorem{remark}{Remark}
\newcommand{\EqRel}{\mathsf{EqRel}}
\newcommand{\Aut}{\operatorname{Aut}}

\title{V21 P1-T04 Proposal-Neutral EqSrc Canonical-Selection\\
Naturality Theorem Target v1}
\author{Ontology Formalizer control packet RT-20260720-012}
\date{2026-07-20}

\begin{document}
\maketitle

\section{Control status and exact burden}

This artifact executes only v21 work item \texttt{P1-T04}. It replaces
candidate-specific selector prose with one common mathematical target for
roots, lines, orientations, partitions, response tokens, normalizations,
ordered actions, components, and analogous distinguished structures. The
target identifier is
\texttt{EQSRC-CANONICAL-SELECTION-NATURALITY-TARGET-V1}.

The record is \texttt{draft/control}, \texttt{proposal-only}, and
\texttt{source-extension data}. Its adoption status is
\texttt{blocked\_adoption\_open\_continuation}. Current ontology does not
derive the source groupoid, the choice-space functor, a natural selector, a
physical morphism category, or a variation law. This packet proves one exact
groupoid lemma below; it does not prove the broader theorem target, construct
a preferred selector, adopt a source law, edit canonical ontology, or
discharge general \(\mathsf{EqSrc}\).

\section{Proposal-neutral source data}

\begin{definition}[Source symmetry domain]
Let \(\mathcal G\) be a small groupoid whose objects are admitted source
records and whose arrows are source isomorphisms declared before any
distinguished choice or candidate relation is evaluated. Let
\[
 A:\mathcal G\longrightarrow\mathsf{Set}
\]
be a carrier functor. For an arrow \(f:S\to T\), transport an equivalence
relation \(R\) on \(A(S)\) to the relation \(f_*R\) on \(A(T)\) by
\[
 y\,(f_*R)\,y'
 \quad\Longleftrightarrow\quad
 A(f)^{-1}(y)\,R\,A(f)^{-1}(y').
\]
Because every arrow of \(\mathcal G\) is invertible, this defines a functor
\(\EqRel(A):\mathcal G\to\mathsf{Set}\).
\end{definition}

\begin{definition}[Choice space, relation map, and selector]
Let \(X:\mathcal G\to\mathsf{Set}\) assign the admissible distinguished
structures. Let
\[
 K:X\Longrightarrow\EqRel(A)
\]
be a natural transformation that maps each choice to its induced candidate
relation. A \emph{natural selector} is a natural transformation
\[
 s:1_{\mathcal G}\Longrightarrow X,
 \qquad X(f)(s_S)=s_T
\]
for every \(f:S\to T\), where \(1_{\mathcal G}\) is the constant singleton
functor. These data are the theorem target, not current-ontology conclusions.
\end{definition}

\section{Exact groupoid component criterion}

Choose one representative \(S_i\) from each connected component of
\(\mathcal G\), and write
\[
 X(S_i)^{\Aut_{\mathcal G}(S_i)}
 =\{x\in X(S_i):X(g)x=x\text{ for every }g\in\Aut_{\mathcal G}(S_i)\}.
\]

\begin{theorem}[Component fixed-point criterion]
There is a bijection
\[
 \operatorname{Nat}(1_{\mathcal G},X)
 \cong
 \prod_i X(S_i)^{\Aut_{\mathcal G}(S_i)}.
\]
Thus a natural selector exists exactly when every component representative
has a nonempty automorphism-fixed locus.
\end{theorem}

\begin{proof}
Given a natural selector \(s\), naturality at an automorphism
\(g:S_i\to S_i\) gives \(X(g)s_{S_i}=s_{S_i}\). Evaluation therefore maps
\(s\) into the displayed product.

Conversely, choose \(x_i\) in each fixed locus. For any object \(S\) in the
component of \(S_i\), choose an isomorphism \(u:S_i\to S\) and set
\(s_S=X(u)x_i\). If \(v:S_i\to S\) is another choice, then
\(v^{-1}u\in\Aut_{\mathcal G}(S_i)\), so
\[
 X(u)x_i=X(v)X(v^{-1}u)x_i=X(v)x_i.
\]
The value is therefore independent of transport. For any \(f:S\to T\),
the composite \(fu:S_i\to T\) is a permitted transport, hence
\(X(f)s_S=X(fu)x_i=s_T\). This defines a natural selector. Evaluation and
transport are inverse constructions.
\end{proof}

\begin{corollary}[No-fixed-point branch]
If one component representative has an empty fixed locus, no natural
selector exists on that component.
\end{corollary}

\begin{corollary}[Relation image branches]
For a component representative \(S_i\), let
\[
 \mathcal R_i=K_{S_i}\bigl(X(S_i)^{\Aut(S_i)}\bigr).
\]
If \(\mathcal R_i\) contains two relations, source symmetry and naturality do
not determine a unique relation. If the fixed locus is nonempty and
\(\mathcal R_i\) is a singleton, the induced relation is unique on the
groupoid component even if the selector is not. The latter is a
choice-irrelevance result, not physical admissibility.
\end{corollary}

\begin{proof}
The theorem transports every fixed choice to a natural selector on its
component, and naturality of \(K\) transports its relation. Distinct elements
of \(\mathcal R_i\) therefore give distinct natural relation sections, while
a singleton image makes every natural selector induce the same relation.
\end{proof}

\section{What remains open}

The component theorem is sufficient only for the declared source groupoid.
For a broader category containing noninvertible source morphisms, a family of
componentwise automorphism-fixed choices need not satisfy all arrow
coherence equations. Independently declared variations add still further
robustness equations. Accordingly, the full theorem target remains open
until a packet:
\begin{enumerate}[label=(O\arabic*)]
  \item derives or obstructs the source domain and \(X\) from registered
  source theory;
  \item states and checks every noninvertible-morphism coherence equation;
  \item proves relation uniqueness or choice irrelevance for the intended
  domain;
  \item defines an independently justified variation family and tests
  robustness; and
  \item audits hidden target, observer, benchmark, generated, registry,
  validator, role, handoff, file-order, and task-process imports.
\end{enumerate}

Supplying a section, root, line, orientation, partition, response token,
normalization, action, component boundary, morphism class, or variation class
without deriving it adds the missing choice data. It does not resolve the
selection burden. Likewise, a natural isomorphism \(h:X\Rightarrow X'\) that
preserves \(K\) transports natural sections and fixed-locus relation images
bijectively; it is a relabeling, not a reduction of the burden.

\section{Finite controls}

\begin{longtable}{p{0.22\linewidth}p{0.68\linewidth}}
\toprule
Control & Exact result \\
\midrule
\endfirsthead
\toprule
Control & Exact result \\
\midrule
\endhead
\texttt{CTRL-POS-UNIQUE-FIXED-001} & A one-object trivial groupoid with
singleton \(X\) has exactly one selector and one induced relation. \\
\texttt{CTRL-NEG-FREE-C2-002} & In the one-object groupoid \(B(C_2)\), let
the nonidentity element swap two choices. The fixed locus is empty, so no
natural selector exists. \\
\texttt{CTRL-NEG-MULTIREL-003} & A one-object trivial groupoid with two
choices and two distinct \(K\)-images has two selectors and relation
nonuniqueness. \\
\texttt{CTRL-POS-IRRELEVANT-004} & The same two-choice groupoid with constant
\(K\) has selector nonuniqueness but relation-level choice irrelevance. \\
\texttt{CTRL-NEG-NONINVERTIBLE-005} & In the category
\(0\xrightarrow{a}1\), take \(X(0)=\{p\}\), \(X(1)=\{q,r\}\), and
\(X(a)(p)=q\). All automorphism actions are trivial, so \(r\) is fixed, but
no natural section can use \(r\) at object 1. Automorphism fixedness alone is
not sufficient beyond groupoids. \\
\bottomrule
\end{longtable}

\section{Prior countermodel map}

\begin{longtable}{p{0.25\linewidth}p{0.16\linewidth}p{0.49\linewidth}}
\toprule
Tracked family & Target branch & Scoped evidence \\
\midrule
\endfirsthead
\toprule
Tracked family & Target branch & Scoped evidence \\
\midrule
\endhead
Intrinsic discriminator & Multiple relation fixed points & Four valid
differentials on fixed carriers induce four relations; current ontology does
not choose the differential. \\
Cycle-boundary line & No fixed point & Forgetting the mark exposes full
\(\mathrm{GL}(2,2)\) symmetry with no invariant proper line. \\
Orientation torsor & Multiple relation fixed points & Three equivariant
descended line families survive conditional algebra while full symmetry
selects none. \\
Rooted partition & No fixed point & Directed-cycle symmetry has no invariant
singleton root, and admitted roots induce different relations. \\
Graded-orbit root & Morphism or variation failure & A freely supplied action
realizes every nontrivial ordered finite partition, while an isolated window
does not establish ambient-component completeness. \\
\bottomrule
\end{longtable}

These are mappings into a shared vocabulary, not new proofs about the full
domain. The five exact candidate families and seven immutable candidates
remain locally frozen. The map does not reject them, reopen them, prove a
global no-go, or preclude a conservative future source extension.

\section{Source purity and Distance-to-GR boundary}

The source records, arrows, carrier functor, choice spaces, relation map, and
variations must be defined without target topology, target atlas, target
metric, proper time, detector semantics, benchmark success, generated
derivatives, registries, validators, roles, handoffs, file order, or task
process. A validator PASS checks only the packet's structural and finite
claims.

The \texttt{source\_equivalence\_eqsrc} milestone remains at
\texttt{draft object exists}. General \(\mathsf{EqSrc}\), \(M_{\rm src}\),
\(g_{\rm eff}\), matter coupling, Einstein equations, exact-GR recovery, and
benchmark promotion remain blocked. The Distance-to-GR and metric-use
ledgers are unchanged. No canonical ontology edit, candidate adoption or
rejection, Gate Chair verdict, publication, or external action is authorized.

The next dependency-ready v21 item is \texttt{P1-T05}, a bounded literature
review of natural operations and no-natural-choice results. It may compare
the registered target with primary sources but may not convert literature
analogy into project proof or adoption.

\end{document}
